\section{Platform overview and report scope}

This report describes the July 2026 version of C-LARA-2 as a current project
snapshot. It draws on the implemented platform, repository source code and
tests, roadmap and issue records, operational documentation, selected
experiments and human collaborator experience. The aim is not to provide a
controlled evaluation of every capability, but to give an inspectable account
of the platform's principal functional areas, the development model through
which they are being extended, and the limitations that remain.

At the highest level, C-LARA-2 supports the construction and use of
multimodal language-learning projects. Users can create or import texts;
annotate them with linguistic, translation and audio information; generate
and select images; compile the results as interactive web resources; and
manage projects and their versions. The platform retains the C-LARA idea that
these activities are connected stages in the production of learner-facing
material, while providing more systematic support for specialised workflows,
revision and operational visibility.

The first group of functions concerns text creation and linguistic
processing. A user can begin with source text or an AI-generated text, choose
relevant project and glossing languages, and run an annotation pipeline that
includes segmentation, translation, MWE analysis, lemma/POS tagging and
glossing. The pipeline is not treated as an infallible black box: users can
review and edit outputs, and the repository retains prompts, tests,
experiments and issue records through which processing behaviour can be
examined and improved.

A second group concerns multimodal presentation. Image-generation workflows
can produce page illustrations and support consistent visual style; picture
dictionaries adapt the underlying project representation to a task in which a
lexical entry, its translation, its image-generation prompt and its selected
image should be managed together. The same resources can support picture
glossing and picture-based exercises. These functions are especially relevant
where accessible visual materials can help learners or community contributors,
but their use remains subject to human pedagogical review and, where
appropriate, cultural constraints.

A third group concerns the infrastructure that makes an evolving platform
usable and inspectable. Project-management features, deployment work,
roadmaps, issue records and operational documentation make it possible to
track work beyond an individual interaction with an AI agent. The
project-understanding Assistant builds on this repository memory: it can
inspect current project artefacts and answer natural-language questions about
implemented behaviour, planned work and design decisions, with links back to
the supporting repository evidence. This is an early, limited form of project
self-understanding rather than a claim of perfect knowledge or autonomous
judgement.

Finally, C-LARA-2 is beginning to explore how prepared multimodal texts might
relate to spoken conversational practice. Sarah Wright's regular sessions with
ChatGPT Pro Tier Voice Mode illustrate a promising conversational complement
to the platform's text-centred resources. The two environments are not yet
naturally integrated, and the report distinguishes present workarounds from
implemented C-LARA-2 functionality.

Sections~4--10 examine these areas in more detail: the annotation pipeline,
image generation, picture dictionaries, project management, deployment, the
project-understanding Assistant and Voice Mode. The organisation is
deliberately selective. It foregrounds functions and episodes that clarify
the C-LARA-2 development model, rather than attempting a complete feature
manual. Sections~11 and~12 then identify the most important unfinished work
and draw together the report's conclusions about operational AI autonomy under
strategic human supervision.
